\src{182}
\[
\symup{\operatorname{Ext}_{\Osh_X}^{p}(X;\Osh_X,F)
\times\operatorname{Ext}_{\Osh_X}^{n-p}(X;F,\underline{\Omega}_X^{n})
\longrightarrow
\operatorname{Ext}_{\Osh_X}^{n}(X;\Osh_X,\underline{\Omega}_X^{n}),\quad\text{i.e.}}
\]
\begin{equation}\tag{5.1}
  \symup{H^{p}(X,F)
  \times\operatorname{Ext}_{\Osh_X}^{n-p}(X;F,\underline{\Omega}_X^{n})
  \longrightarrow H^{n}(X,\underline{\Omega}_X^{n})}
\end{equation}
Compte tenu du théorème 2, cet accouplement définit un homomorphisme
\begin{equation}\tag{5.2}
  \symup{\operatorname{Ext}_{\Osh_X}^{n-p}(X;F,\underline{\Omega}_X^{n})
  \longrightarrow\bigl(H^{p}(X,F)\bigr)'}
\end{equation}
Cet homomorphisme est fonctoriel en $\symup{F}$, et permute aux homomorphismes cobords
relatifs aux suites exactes $\symup{0\longrightarrow F'\longrightarrow F
\longrightarrow F''\longrightarrow0}$.

\result{THÉORÈME}{3} \uline{L'homomorphisme} (5.2) \uline{est un isomorphisme}.

On retrouve en particulier le résultat de Serre :

\result{COROLLAIRE}{} \uline{Soient} $\symup{E}$ \uline{un fibré vectoriel algébrique sur} $\symup{X}$,
$\symup{\Osh_X(E)}$ \uline{le faisceau des germes de sections régulières de} $\symup{E}$, \uline{on a alors des
isomorphismes canoniques} :
\begin{equation}\tag{5.3}
  \symup{\bigl(H^{p}(X,\Osh_X(E))\bigr)'
  =H^{n-p}\bigl(X,\underline{\Omega}_X^{n}\otimes\Osh_X(E')\bigr)}
\end{equation}
Il suffit d'appliquer le théorème 3 et le corollaire 1 à la proposition 1.

Les théorèmes 2 et 3 vont résulter de l'énoncé suivant :

\noindent(D) \uline{L'homomorphisme suivant}
\begin{equation}\tag{5.2 bis}
  \symup{\operatorname{Ext}_{\Osh_X}^{n-p}(X;F,\underline{\Omega}_X^{n})
  \longrightarrow\bigl(H^{p}(X,F)\bigr)'\otimes L
  \qquad (L=H^{n}(X,\underline{\Omega}_X^{n})),}
\end{equation}
\uline{déduit de l'accouplement (5.1) est un isomorphisme}.

Montrons en effet que (D) implique le théorème 2. Soit
$\symup{\underline{k}_x=\Osh_{(x)}}$ le faisceau des anneaux locaux de la variété réduite au point
$\symup{x\in X}$, considérons l'homomorphisme canonique
$\symup{\Osh_X\longrightarrow \underline{k}_x}$ et l'homomorphisme associé
\begin{equation}\tag{5.3 bis}
  \symup{H^{0}(X,\Osh_X)\longrightarrow H^{0}(X,\underline{k}_x)}
\end{equation}
son transposé s'identifie à l'homomorphisme
\begin{equation}\tag{5.4}
  \symup{\operatorname{Ext}_{\Osh_X}^{n}(X;\underline{k}_x,\underline{\Omega}_X^{n})\otimes L'
  \longrightarrow
  \operatorname{Ext}_{\Osh_X}^{n}(X;\Osh_X,\underline{\Omega}_X^{n})\otimes L'}
\end{equation}
déduit de l'homomorphisme sur les $\symup{\operatorname{Ext}^{n}}$ associé à
$\symup{\Osh_X\longrightarrow \underline{k}_x}$ :
\begin{equation}\tag{5.5}
  \symup{\operatorname{Ext}_{\Osh_X}^{n}(X;\underline{k}_x,\underline{\Omega}_X^{n})
  \longrightarrow
  \operatorname{Ext}_{\Osh_X}^{n}(X;\Osh_X,\underline{\Omega}_X^{n})
  =H^{n}(X,\underline{\Omega}_X^{n})}
\end{equation}

